\chapter{Un \texttt{.jenga} est du Python}
\label{ch:python}

Le chapitre 1 l'a affirmé ; celui-ci en tire les conséquences pratiques. Ce
n'est pas une curiosité d'implémentation : c'est ce qui vous permet d'écrire une
boucle là où d'autres systèmes vous demanderaient un langage de macros.

\section{Le squelette obligatoire}

\begin{nkcode}{Le minimum viable}
from Jenga import *

with workspace("MonEspace"):
    configurations(['Debug', 'Release'])

    with project("MonProgramme"):
        consoleapp()
        language("C++")
        files(["src/**.cpp"])
\end{nkcode}

\begin{nkretenir}
\textbf{\texttt{with}, et pas autre chose.} \texttt{workspace} et
\texttt{project} sont des \emph{gestionnaires de contexte} Python --- des classes
avec \texttt{\_\_enter\_\_} et \texttt{\_\_exit\_\_}.

Ce que cela signifie concrètement : \textbf{tout appel écrit à l'intérieur du
bloc s'applique à ce bloc.} \texttt{files()} ne prend pas de projet en
paramètre ; il sait où il est parce qu'il lit une variable globale que
\texttt{with project(\dots)} vient de poser.
\end{nkretenir}

\begin{nkpiege}
\textbf{Un appel écrit HORS d'un bloc ne disparaît pas : il s'applique
ailleurs.} \texttt{files()} après la fin d'un \texttt{with project} s'applique au
projet \emph{précédent}, ou à rien --- pas d'erreur, pas d'avertissement.

L'indentation Python n'est donc pas cosmétique ici : \textbf{elle décide à quoi
un réglage appartient}. Un décalage d'un niveau change la signification sans
jamais échouer.
\end{nkpiege}

\section{Ce que le Python vous donne, et qu'aucun format ne donnerait}

\subsection{Une boucle}

\begin{nkcode}{Six demos, une description}
for i in range(1, 7):
    with project("Demo%d" % i):
        consoleapp()
        language("C++")
        location("Demos/Demo%d" % i)
        files(["src/**.cpp"])
\end{nkcode}

\subsection{Une condition sur l'environnement}

Le dépôt s'en sert réellement, pour ne déclarer Vulkan que là où il est câblé :

\begin{nkcode}{Applications/NkDames/NkDames.jenga}
_VULKAN_SDK = _norm(os.getenv("VULKAN_SDK", ""))
_VULKAN_OPT = os.getenv("NK_ENABLE_VULKAN", "auto").strip().lower()
if _VULKAN_OPT in ("0", "false", "off", "no"):
    _WANT_VULKAN = False
elif _VULKAN_OPT in ("1", "true", "on", "yes"):
    _WANT_VULKAN = True
else:
    _WANT_VULKAN = bool(_VULKAN_INCLUDE) or _pkg_exists("vulkan")
\end{nkcode}

\subsection{Une fonction à vous}

\begin{nkcode}{Applications/NkDames/NkDames.jenga}
def _pkg_exists(pkg_name):
    if not shutil.which("pkg-config"):
        return False
    try:
        return subprocess.call(["pkg-config", "--exists", pkg_name],
                               stdout=subprocess.DEVNULL,
                               stderr=subprocess.DEVNULL) == 0
    except Exception:
        return False
\end{nkcode}

\begin{nkretenir}
\textbf{Ce descripteur interroge le système pendant qu'il se charge.} Il lance
\texttt{pkg-config}, lit des variables d'environnement, teste des chemins.

Aucun format déclaratif ne permet cela --- et c'est exactement ce dont on a
besoin quand une bibliothèque est présente sur une machine et pas sur une autre.
\end{nkretenir}

\begin{nkpiege}[Le revers, et il est sérieux]
\textbf{Une description qui interroge le système donne un résultat qui dépend du
système.} Deux machines, deux descriptions effectives, et le
\texttt{.jenga} est pourtant identique dans le dépôt.

C'est puissant et c'est un piège : « ça construit chez moi » devient une phrase
qui ne prouve rien. \textbf{Quand vous écrivez une condition, faites-la dire
quelque chose} --- une ligne de journal --- sinon la différence entre deux
machines reste invisible jusqu'à l'édition de liens.
\end{nkpiege}

\section{Les erreurs sont des erreurs Python}

\begin{nkretenir}
Le chargeur attrape et affiche :

\begin{nkcode}{Jenga/Core/Loader.py}
except Exception as e:
    Colored.PrintError(f"Error loading workspace: {e}")
    if self.verbose:
        traceback.print_exc()
\end{nkcode}

\textbf{Avec \texttt{-{}-verbose}, vous obtenez la pile d'appels complète.} Sans
lui, vous n'avez que le message --- souvent suffisant, parfois non. C'est le
premier réflexe quand un \texttt{.jenga} refuse de se charger.
\end{nkretenir}

\begin{center}
\begin{tabular}{@{}ll@{}}
\toprule
\textbf{Ce que vous voyez} & \textbf{Ce que c'est} \\
\midrule
\texttt{SyntaxError} & une virgule, une parenthèse, une indentation \\
\texttt{NameError} & une fonction Jenga mal orthographiée \\
\texttt{TypeError} & un argument du mauvais type --- souvent une chaîne au lieu d'une liste \\
\texttt{FileNotFoundError} & un \texttt{include()} vers un fichier absent \\
\texttt{No workspace defined} & aucun \texttt{with workspace(\dots)} n'a été exécuté \\
\bottomrule
\end{tabular}
\end{center}

\begin{nkpiege}
\textbf{La plupart des fonctions attendent une LISTE.} \texttt{files("src/main.cpp")}
--- sans crochets --- ne donne pas toujours une erreur : une chaîne est itérable,
et Python la parcourt \emph{caractère par caractère}.

Vous obtenez alors un projet dont les fichiers s'appellent \texttt{s},
\texttt{r}, \texttt{c}, \texttt{/}\dots\ et le message final parlera de fichiers
introuvables, sans jamais nommer la cause.
\end{nkpiege}

\section{Le répertoire courant change pendant le chargement}

\begin{nkcode}{Jenga/Core/Loader.py}
old_cwd = Path.cwd()
os.chdir(filePath.parent)
\end{nkcode}

\begin{nkretenir}
\textbf{Pendant l'exécution d'un \texttt{.jenga}, le répertoire courant est celui
du fichier.} C'est ce qui rend les chemins relatifs prévisibles : ils sont
relatifs au descripteur, pas à l'endroit d'où vous avez tapé la commande.

Et c'est vrai \emph{récursivement} : un fichier inclus s'exécute dans
\emph{son} dossier. Un \texttt{include("Applications/X/X.jenga")} suivi de
\texttt{location(".")} dans le fichier inclus désigne bien
\texttt{Applications/X}.
\end{nkretenir}

\section{\texttt{include} : un seul workspace, plusieurs fichiers}

\begin{nkcode}{Nkentseu.jenga, a la racine}
with include("Applications/NkDames/NkDames.jenga"):
    pass
\end{nkcode}

\begin{nkretenir}
\textbf{Le fichier inclus s'exécute \emph{en ligne}, dans le workspace courant.}
Il ne crée pas un second workspace : il ajoute ses projets au vôtre.

\textbf{Et il hérite du \emph{namespace} du fichier racine.} Le chargeur expose
les variables globales du fichier d'entrée aux inclusions :

\begin{nkcode}{Jenga/Core/Loader.py}
# Exposer le namespace LIVE du workspace racine pour que `include()` propage
# automatiquement aux fichiers inclus tout ce que l'utilisateur definit ici
# (config/fonctions/classes/imports) -> plus besoin de les re-importer dans
# chaque .jenga inclus.
\end{nkcode}

Une fonction ou une constante définie \emph{au-dessus} des inclusions est donc
visible dans tous les fichiers inclus. \textbf{Au-dessus} : les inclusions
s'exécutent à l'endroit où elles sont écrites, et ne voient que ce qui les
précède.
\end{nkretenir}

\begin{nkpiege}
\textbf{Un projet non inclus n'existe pas.} \texttt{jenga build --target MonJeu}
répond que la cible est inconnue, et l'on relit son \texttt{.jenga} vingt fois ---
qui est parfaitement juste.

La question à se poser n'est pas « qu'ai-je mal écrit dans ce fichier ? » mais
« ce fichier est-il inclus ? ». \texttt{jenga info} tranche en une seconde : le
projet est dans le tableau, ou il n'y est pas.
\end{nkpiege}

\begin{nkbilan}
Un \texttt{.jenga} est du Python exécuté, et \texttt{workspace} comme
\texttt{project} sont des gestionnaires de contexte : l'indentation décide à
quoi un réglage appartient, et un appel mal placé s'applique ailleurs sans rien
dire. Vous disposez des boucles, des conditions et de vos propres fonctions ---
au prix qu'une description peut alors dépendre de la machine, ce qu'il faut
rendre visible. Les erreurs sont des erreurs Python, et \texttt{-{}-verbose} en
donne la pile. Pendant le chargement, le répertoire courant est celui du fichier,
récursivement. Enfin \texttt{include} exécute en ligne, dans le même workspace,
en propageant le namespace de ce qui précède.
\end{nkbilan}

\begin{nkquestions}
\begin{enumerate}
\item Pourquoi \texttt{files()} n'a-t-il pas besoin qu'on lui dise de quel projet
      il s'agit ?
\item Que se passe-t-il si un appel est écrit hors de son bloc ?
\item Que produit \texttt{files("src/main.cpp")} sans crochets, et pourquoi le
      message d'erreur ne le dit-il pas ?
\item Quel est le répertoire courant pendant l'exécution d'un \texttt{.jenga} ?
\item Une fonction définie \emph{après} un \texttt{include} est-elle visible dans
      le fichier inclus ? Justifiez.
\end{enumerate}
\end{nkquestions}

\begin{nkpratique}
Écrivez un workspace qui déclare \textbf{cinq} projets identiques à la casse du
nom près, \emph{sans copier-coller} : une boucle.

Puis ajoutez une condition : si la variable d'environnement \texttt{DEMO\_LEGER}
vaut \texttt{1}, seuls les deux premiers sont déclarés.

Vérifiez par \texttt{jenga info} dans les deux cas --- \textbf{sans construire}.
C'est le contrôle le moins cher, et il suffit ici.
\end{nkpratique}

\begin{nkdemo}
Prouvez que l'indentation change la signification, et pas seulement la
présentation.

Écrivez deux projets. Dans le premier essai, placez un \texttt{defines(["A"])}
\emph{à l'intérieur} du second bloc ; dans le second essai, à l'extérieur, aligné
sur les \texttt{with project}.

Comparez les deux par \texttt{jenga compile-flags} --- qui affiche les drapeaux
réels de chaque projet.

\textbf{Ce que la démonstration doit établir} : dans un cas la définition est
attachée à un projet, dans l'autre elle a changé de propriétaire ou disparu ---
\textbf{sans qu'aucune des deux exécutions n'ait échoué}. C'est cette absence
d'erreur qui rend le défaut coûteux.
\end{nkdemo}

\begin{nklogique}
Un descripteur construit sur votre machine et échoue sur celle d'un collègue,
avec une erreur d'édition de liens qui nomme un symbole Vulkan.

\begin{enumerate}
\item Le fichier étant identique dans le dépôt, expliquez comment les deux
      descriptions peuvent différer.
\item Où, dans le fichier, chercheriez-vous en premier ? Nommez la construction
      Python responsable.
\item Proposez une modification qui rende la différence \emph{visible} au
      chargement, avant l'édition de liens --- et dites pourquoi la placer
      ailleurs qu'à cet endroit ne servirait à rien.
\end{enumerate}
\end{nklogique}
